
\begin{table}[htbp]
   
   \caption{\label{tab:fs_dwnstrm_nonminevio_ivsum} First stage: effect of the Acid Rain Program on the number of upstream coal mines (summed across upstream watersheds, utilities at most one HUC12 downstream)}
   \bigskip
   
   \centering
   
   \begin{adjustbox}{width = 0.45\textwidth, center}
      
      %start:tab
      
      \begin{tabular}{lrr}
         \toprule
         
          & \multicolumn{2}{c}{Upstream coal mines (sum)}\\
         
          & \multicolumn{2}{c}{any violation category} \\ 
                                              & Total coliform  & VOCs \\   
         
                                              & (1)             & (2)\\  
         
         \midrule 
         post95 $\times$ Upstream sulfur \%   & -0.1469$^{***}$ & -0.1763$^{***}$\\   
                                              & (0.0354)        & (0.0384)\\   
         
         F-test (1st stage, clustered)        & 17.21           & 21.07\\  
         
          \\
         
         Observations                         & 4,254           & 4,576\\  
         
         \bottomrule
         
      \end{tabular}
      
      %end:tab
      
      
   \end{adjustbox}
      {\tiny\linespread{1}\selectfont \par \raggedright 
   \textit{Notes:} The dependent variable is the number of coal mines in watersheds upstream of the utility's intake, summed across upstream watersheds. The instrument interacts an indicator for the post-1995 period with the sum of coal sulfur content across upstream watersheds. The sample is utilities at most one watershed downstream of a coal mine. All specifications include utility and year fixed effects. Standard errors clustered at the utility level. Sample period 1985--2005. *** p$<$0.01, ** p$<$0.05, * p$<$0.1.}
   
\end{table}


